\chapter{2018年北京卷（理）}

\section{单选题}

\begin{problem}
已知集合 \(A=\{x\mid |x|<2\}\)，\(B=\{-2,0,1,2\}\)，则 \(A\cap B=\)（\quad）.

\choices
    {\(\{0,1\}\)}
    {\(\{-1,0,1\}\)}
    {\(\{-2,0,1,2\}\)}
    {\(\{-1,0,1,2\}\)}

\end{problem}
\begin{problem}
在复平面内，复数 \(\dfrac{1}{1-\i}\) 的共轭复数对应的点位于（\quad）.

\choices
    {第一象限}
    {第二象限}
    {第三象限}
    {第四象限}

\end{problem}

\begin{problem}
执行如图所示的程序框图，输出的 \(s\) 值为（\quad）.

\bitmapfigure[max width=0.32\linewidth]{img/2018/beijing_science/q3_fig1.png}

\choices
    {\(\dfrac12\)}
    {\(\dfrac56\)}
    {\(\dfrac76\)}
    {\(\dfrac{7}{12}\)}

\end{problem}

\begin{problem}
“十二平均律”是通用的音律体系，明代朱载堉最早用数学方法计算出半音比例，为这个理论的发展做出了重要贡献. 十二平均律将一个纯八度音程分成十二份，依次得到十三个单音，从第二个单音起，每一个单音的频率与它的前一个单音的频率的比都等于 \(2^{\frac1{12}}\). 若第一个单音的频率为 \(f\)，则第八个单音的频率为（\quad）.

\choices
    {\(\sqrt[3]{2}\,f\)}
    {\(\sqrt[3]{2^2}\,f\)}
    {\(\sqrt[12]{2^5}\,f\)}
    {\(\sqrt[12]{2^7}\,f\)}

\end{problem}

\begin{problem}
某四棱锥的三视图如图所示，在此四棱锥的侧面中，直角三角形的个数为（\quad）.

\bitmapfigure[max width=6.0cm]{img/2018/beijing_science/q5_fig1.png}

\choices
    {\(1\)}
    {\(2\)}
    {\(3\)}
    {\(4\)}

\end{problem}

\begin{problem}
设 \(\bs a\)，\(\bs b\) 均为单位向量，则“\(\left\lvert \bs a-3\bs b\right\rvert=\left\lvert 3\bs a+\bs b\right\rvert\)”是“\(\bs a\perp \bs b\)”的（\quad）.

\choices
    {充分而不必要条件}
    {必要而不充分条件}
    {充分必要条件}
    {既不充分也不必要条件}

\end{problem}

\begin{problem}
在平面直角坐标系中，记 \(d\) 为点 \(P(\cos\theta,\sin\theta)\) 到直线 \(x-my-2=0\) 的距离，当 \(\theta\)，\(m\) 变化时，\(d\) 的最大值为（\quad）.

\choices
    {\(1\)}
    {\(2\)}
    {\(3\)}
    {\(4\)}

\end{problem}

\begin{problem}
设集合 \(A=\{(x,y)\mid x-y\ge1,ax+y>4,x-ay\le2\}\)，则（\quad）.

\choices
    {对任意实数 \(a\)，\((2,1)\in A\)}
    {对任意实数 \(a\)，\((2,1)\notin A\)}
    {当且仅当 \(a<0\) 时，\((2,1)\notin A\)}
    {当且仅当 \(a\le \dfrac32\) 时，\((2,1)\notin A\)}

\end{problem}

\section{填空题}

\begin{problem}
设 \(\{a_n\}\) 是等差数列，且 \(a_1=3\)，\(a_2+a_5=36\)，则 \(\{a_n\}\) 的通项公式为\fillinblank{}.

\end{problem}

\begin{problem}
在极坐标系中，直线 \(\rho\cos\theta+\rho\sin\theta=a\)（\(a>0\)）与圆 \(\rho=2\cos\theta\) 相切，则 \(a=\fillinblank{}\).

\end{problem}

\begin{problem}
设函数 \(f(x)=\cos\!\left(\omega x-\dfrac{\pi}{6}\right)\)（\(\omega>0\)），若 \(f(x)\le f\!\left(\dfrac{\pi}{4}\right)\) 对任意实数 \(x\) 都成立，则 \(\omega\) 的最小值为\fillinblank{}.

\end{problem}

\begin{problem}
若 \(x\)，\(y\) 满足 \(x+1\le y\le 2x\)，则 \(2y-x\) 的最小值是\fillinblank{}.

\end{problem}

\begin{problem}
能说明“若 \(f(x)>f(0)\) 对任意的 \(x\in(0,2]\) 都成立，则 \(f(x)\) 在 \([0,2]\) 上是增函数”为假命题的一个函数是\fillinblank{}.

\end{problem}

\begin{problem}
已知椭圆 \(M:\dfrac{x^2}{a^2}+\dfrac{y^2}{b^2}=1\)（\(a>b>0\)），双曲线 \(N:\dfrac{x^2}{m^2}-\dfrac{y^2}{n^2}=1\). 若双曲线 \(N\) 的两条渐近线与椭圆 \(M\) 的四个交点及椭圆 \(M\) 的两个焦点恰为一个正六边形的顶点，则椭圆 \(M\) 的离心率为\fillinblank{}；双曲线 \(N\) 的离心率为\fillinblank{}.

\end{problem}

\section{解答题}

\begin{problem}
在 \(\triangle ABC\) 中，\(a=7\)，\(b=8\)，\(\cos B=-\dfrac17\).

\begin{enumerate}
    \item 求 \(\angle A\)；
    \item 求 \(AC\) 边上的高.
\end{enumerate}

\end{problem}

\begin{problem}
如图，在三棱柱 \(ABC-A_1B_1C_1\) 中，\(CC_1\perp\) 平面 \(ABC\)，\(D\)，\(E\)，\(F\)，\(G\) 分别为 \(AA_1\)，\(AC\)，\(A_1C_1\)，\(BB_1\) 的中点，\(AB=BC=\sqrt5\)，\(AC=AA_1=2\).

\begin{enumerate}
    \item 求证：\(AC\perp\) 平面 \(BEF\)；
    \item 求二面角 \(B-CD-C_1\) 的余弦值；
    \item 证明：直线 \(FG\) 与平面 \(BCD\) 相交.
\end{enumerate}

\bitmapfigure[max width=0.62\linewidth]{img/2018/beijing_science/q16_fig1.png}

\end{problem}

\begin{problem}
电影公司随机收集了电影的有关数据，经分类整理得到下表：

\begin{center}
\begin{tabular}{c|cccccc}
\hline
电影类型 & 第一类 & 第二类 & 第三类 & 第四类 & 第五类 & 第六类 \\
\hline
电影部数 & 140 & 50 & 300 & 200 & 800 & 510 \\
好评率 & 0.4 & 0.2 & 0.15 & 0.25 & 0.2 & 0.1 \\
\hline
\end{tabular}
\end{center}

好评率是指：一类电影中获得好评的部数与该类电影的部数的比值.

假设所有电影是否获得好评相互独立.

\begin{enumerate}
    \item 从电影公司收集的电影中随机选取 \(1\) 部，求这部电影是获得好评的第四类电影的概率；
    \item 从第四类电影和第五类电影中各随机选取 \(1\) 部，估计恰有 \(1\) 部获得好评的概率；
    \item 假设每类电影得到人们喜欢的概率与表格中该类电影的好评率相等，用“\(\xi_k=1\)”表示第 \(k\) 类电影得到人们喜欢，“\(\xi_k=0\)”表示第 \(k\) 类电影没有得到人们喜欢（\(k=1\)，\(2\)，\(3\)，\(4\)，\(5\)，\(6\)）. 写出方差 \(D(\xi_1)\)，\(D(\xi_2)\)，\(D(\xi_3)\)，\(D(\xi_4)\)，\(D(\xi_5)\)，\(D(\xi_6)\) 的大小关系.
\end{enumerate}

\end{problem}

\begin{problem}
设函数 \(f(x)=\bigl[ax^2-(4a+1)x+4a+3\bigr]\e^x\).

\begin{enumerate}
    \item 若曲线 \(y=f(x)\) 在点 \((1,f(1))\) 处的切线与 \(x\) 轴平行，求 \(a\)；
    \item 若 \(f(x)\) 在 \(x=2\) 处取得极小值，求 \(a\) 的取值范围.
\end{enumerate}

\end{problem}

\begin{problem}
已知抛物线 \(C:y^2=2px\) 经过点 \(P(1,2)\). 过点 \(Q(0,1)\) 的直线 \(l\) 与抛物线 \(C\) 有两个不同的交点 \(A\)，\(B\)，且直线 \(PA\) 交 \(y\) 轴于 \(M\)，直线 \(PB\) 交 \(y\) 轴于 \(N\).

\begin{enumerate}
    \item 求直线 \(l\) 的斜率的取值范围；
    \item 设 \(O\) 为原点，\(QM=\lambda QO\)，\(QN=\mu QO\)，求证：\(\dfrac1\lambda+\dfrac1\mu\) 为定值.
\end{enumerate}

\end{problem}

\begin{problem}
设 \(n\) 为正整数，集合 \(A=\left\{(t_1,t_2,\ldots,t_n)\mid t_i\in\{0,1\},\ i=1,2,\ldots,n\right\}\).
对于集合 \(A\) 中的任意元素 \(\alpha=(x_1,x_2,\ldots,x_n)\) 和 \(\beta=(y_1,y_2,\ldots,y_n)\)，记 \(M(\alpha,\beta)=\frac12\sum_{i=1}^n\bigl(x_i+y_i-|x_i-y_i|\bigr)\).

\begin{enumerate}
    \item 当 \(n=3\) 时，若 \(\alpha=(1,1,0)\)，\(\beta=(0,1,1)\)，求 \(M(\alpha,\alpha)\) 和 \(M(\alpha,\beta)\) 的值；
    \item 当 \(n=4\) 时，设 \(B\) 是 \(A\) 的子集，且满足：对于 \(B\) 中的任意元素 \(\alpha\)，\(\beta\)，当 \(\alpha\)，\(\beta\) 相同时，\(M(\alpha,\beta)\) 是奇数；当 \(\alpha\)，\(\beta\) 不同时，\(M(\alpha,\beta)\) 是偶数. 求集合 \(B\) 中元素个数的最大值；
    \item 给定不小于 \(2\) 的 \(n\)，设 \(B\) 是 \(A\) 的子集，且满足：对于 \(B\) 中的任意两个不同的元素 \(\alpha\)，\(\beta\)，\(M(\alpha,\beta)=0\). 写出一个集合 \(B\)，使其元素个数最多，并说明理由.
\end{enumerate}

\end{problem}
